\documentclass[10 pt, a4paper, twoside]{report}

\usepackage{packages}

\makenoidxglossaries
\loadglsentries{gloss}

%\usepackage{newcommands}
\newcommand{\corch}[1]{
\left\lbrace #1 \right\rbrace
}

\newcommand{\sucn}[1]{
\lbrace #1_n \rbrace_{n=1}^{\infty}
}

\newcommand{\suci}[1]{
\lbrace #1_i \rbrace_{i=1}^{\infty}
}

\newcommand{\cupn}[1]{
\cup_{n=1}^\infty #1_n 
}

\newcommand{\sumn}[1]{
\sum_{n=1}^\infty #1
}


\newcommand{\parent}[1]{
\left( #1 \right)
}

\newcommand{\abs}[1]{
\left| #1 \right|
}


\newcommand{\Definition}[2]{
\Bluebox{\begin{definition}[#1]
#2
\end{definition}}
}

\newcommand{\Theorem}[2]{
\Yellowbox{\begin{theorem}[#1]
#2
\end{theorem}}
}

\newcommand{\Proposition}[2]{
\Redbox{\begin{prop}[#1]
#2
\end{prop}}
}

\newcommand{\Lemma}[2]{
\Greenbox{\begin{lemma}[#1]
#2
\end{lemma}}
}

\newcommand{\R}[0]{
\mathbb{R}
}

\newcommand{\K}[0]{
\mathbb{K}
}

\newcommand{\Z}[0]{
\mathbb{Z}
}

\newcommand{\C}[0]{
\mathbb{C}
}

\newcommand{\Q}[0]{
\mathbb{Q}
}

\newcommand{\N}[0]{
\mathbb{N}
}


\newcommand{\CP}[0]{
\mathbb{CP}
}

\newcommand{\Proy}{
    \mathbb{P}
}

\newcommand{\Proj}{
    \operatorname{Proj}
}

\renewcommand{\S}{
    \mathbb{S}
}

\newcommand{\A}{
    \mathbb{A}
}

\newcommand{\norm}[1]{
\lVert #1 \rVert
}

\newcommand{\zenterg}{
    \mathfrak{z}(\mathfrak{g})
}

\newcommand{\Spec}{
    \operatorname{Spec}
}

\newcommand{\diff}[0]{
\mathrm{d}
}

\newcommand{\ii}[0]{
    \mathrm{i} \hspace{2 pt}
}

\newcommand{\Ree}[1]{
    \mathrm{Re}(#1) \hspace{2 pt}
}

\newcommand{\Imm}[1]{
    \mathrm{Im}(#1) \hspace{2 pt}
}

\newcommand{\basecoord}[2]{
\frac{\partial}{\partial #1} \bigg|_{#2}
}

\newcommand{\barpartial}[0]{
\overline{\partial}
}

\newcommand{\basevect}[1]{
\frac{\partial}{\partial #1} 
}

\newcommand{\derpar}[2]{
    \frac{\partial #1}{\partial #2}
}

\newcommand{\Ad}[1]{
    \mathrm{Ad}_{#1}
}

\newcommand{\Diff}[0]{
    \mathrm{D}
}


\newcommand{\bigdt}[0]{
    \left. \frac{\diff}{\diff t} \right|_{t=0}
}

\theoremstyle{definition}
\newtheorem{theorem}{Theorem}[chapter]
\newtheorem{definition}[theorem]{Definition}
\newtheorem{corollary}[theorem]{Corollary}
\newtheorem{prop}[theorem]{Proposition}
\newtheorem{lemma}[theorem]{Lemma}
\newtheorem{example}[theorem]{Example}
\newtheorem{obs}[theorem]{Remark}
\newtheorem{nota}[theorem]{Notation}

%-----------------------------------------------------------------------------------


% Esto es un paquete que deja poner índices
% \usepackage{imakeidx}
% \makeindex[title=Subject index, intoc]


% MUY IMPORTANTE SI SE QUIERE COMPILAR UN DOCUMENTO CON ÍNDICE ALFABÉTICO AL FINAL
% hay que escribir en la terminal
% makeindex main.idx -s mystyle.ist
% Aquí main.idx es el archivo que se va a crear con las entradas que van a salir en el índice

\usepackage{makeidx}
\makeindex
\makeatletter % Permite usar comandos de LaTeX con @ en su nombre
\renewenvironment{theindex}
  {\begin{multicols}{2}[\chapter*{Subject index}]%
   \thispagestyle{plain}%
   \raggedright%
   \columnsep=35pt%
   \columnseprule=0pt%
   \let\item\@idxitem%  <-- Este comando DEBE estar aquí.
   \parindent=0pt%
   \parskip=0pt \@plus .3\p@\relax%
  }%
  {\end{multicols}%
  % \clearpage%
  }
\makeatother % Restaura el comportamiento normal para @

% para poner hipervínculos
\usepackage{hyperref}
\hypersetup{
%hyperindex = false,
colorlinks = true,
linkcolor = blue,
urlcolor = blue,
citecolor = RoyalBlue,
pdftitle = {Moment map equations in Kähler geometry and gauge theory}
}


%\usepackage[totoc]{idxlayout}
%-----------------------------------------------------------------------------------

% para poner el símbolo de cojunto vacío más bonito
\let\oldemptyset\emptyset
\let\emptyset\varnothing
%---------------------------------------------------------------------------------

\title{TFM}
\author{Sergio Martín Nieto}
\date{June 2026}

\begin{document}

\begin{titlepage}

\vspace{2 cm}
{\huge UNIVERSIDAD COMPLUTENSE DE MADRID \par}
    \vspace{30 pt}
    \begin{figure}[h!]
    \centering
        \includegraphics[scale=0.3]{EscudoUCM.png}
    \end{figure}
\begin{center}    
{\bfseries\LARGE FACULTAD DE CIENCIAS MATEMÁTICAS \par}
    \vspace{30 pt}
{\Large \scshape Departamento de álgebra, geometría y topología \par}
    \vspace{3 cm}
{\scshape\Large TRABAJO DE FIN DE MÁSTER\par}
{\scshape\Large  \par}
\vspace{0.75cm}
{ \Large\textbf{Moment map equations in Kähler geometry and gauge theory} \par}
\vspace{2 cm}
{ Author: Sergio Martín Nieto \par}
{ Supervisor: Óscar García Prada \par}
{ Tutor UCM: Marco Castrillón López \par}
\vspace{0.5 cm}
{ Madrid, 24 de junio de 2026 \par}
    
\end{center}
\end{titlepage}
\(\hspace{1 pt}\)
\thispagestyle{empty}
\newpage

% \thispagestyle{empty}
% $\hspace{1 pt}$
% \newpage
% \thispagestyle{empty}
% $\hspace{1 pt}$ \newline \newline \newline \newline \newline \newline
% \begin{flushright}
% \textit{cita \newline \newline
% \textbf{cita}}
% \end{flushright}
% \newpage
% \thispagestyle{empty}
% $\hspace{1 pt}$
% \newpage


% ABSTRACT
\thispagestyle{empty}
\topskip0pt
%\vspace*{\fill}
\begin{center}
    {\scshape Resumen\par}
\end{center}
    Resumen breve.\newline

    \noindent \textbf{Palabras clave:} 
    \newline \newline \newline \newline

\begin{center}
    {\scshape Abstract\par}
\end{center}
    Brief summary \newline

    \noindent \textbf{Keywords:} 

\vspace*{\fill}

\(\hspace{1pt}\)
\newpage
\(\hspace{1pt}\)
\thispagestyle{empty}
\newpage


{\hypersetup{hidelinks}
\tableofcontents
}

\glsaddall

\newpage
\include*{Introduction} 
\include*{Chapter1}
\include*{Chapter2}
\include*{Chapter3}
% \include*{Chapter4} Todo el contenido de este capítulo está redistribuido, no hace falta compilarlo
\include*{Chapter5}
\include*{ChapterHK}

\appendixpageoff
\begin{appendices}
    \include*{Chapter_kahler}
    \include*{appendix_lie_theory}
    \include*{appendix_Banach_manifolds}
\end{appendices}

%\printindex
%\clearpage

% \markboth{Subject index}{Subject index}%
% \printindex
% \addcontentsline{toc}{chapter}{Subject index}
% \vspace{2em}  % small visual gap instead of a new page
% \printnoidxglossary[type=symbols, sort=use, title={Notation index}]
% \addcontentsline{toc}{chapter}{Notation index}
% \printbibliography[heading=bibintoc, title = References]

{
  \let\clearpage\relax
  \newpage
  \markboth{Subject index}{Subject index}%
  \printindex
  \addcontentsline{toc}{chapter}{Subject index}
  \newpage
  \printnoidxglossary[type=symbols, sort=use, title={Notation index}]
  \addcontentsline{toc}{chapter}{Notation index}
}


\printbibliography[heading=bibintoc, title = References]



% \begingroup
% \glsnogroupskiptrue
% \renewcommand{\glossarysection}[2][]{} % suppress internal TOC+title call
% \newpage
% \begin{center}
%     {\bfseries\scshape\Huge Notation index}
% \end{center}
% \printnoidxglossary[type=symbols, sort=use, title={Notation index}]
% \endgroup
% \addcontentsline{toc}{chapter}{Notation index}

% \defbibheading{bibintoc}{%
%   \phantomsection
%   \chapter*{References}%
%   \markboth{References}{References}%
%   \addcontentsline{toc}{chapter}{References}
% }
% \printbibliography[heading=bibintoc]

%\printnoidxglossary[type=symbols, sort=use, title={Notation index}]
%\printindex

\end{document}
